\documentclass[11pt]{article}

\usepackage[margin=1in]{geometry}
\usepackage{amsmath,amssymb,amsthm,mathtools}
\usepackage{enumitem}
\usepackage{graphicx}
\usepackage{xcolor}
\usepackage{parskip}
\usepackage{hyperref}
\hypersetup{colorlinks=true,linkcolor=blue!45!black,urlcolor=blue!45!black}

%% ------------------------------------------------------------------
%%  Theorem-like environments
%%    plain     : theorem lemma proposition corollary claim
%%    definition: definition example exercise
%%    remark    : remark note   (unnumbered)
%%  Proof environment comes from amsthm: \begin{proof} ... \end{proof}
%% ------------------------------------------------------------------
\theoremstyle{plain}
\newtheorem{theorem}{Theorem}[section]
\newtheorem{lemma}[theorem]{Lemma}
\newtheorem{proposition}[theorem]{Proposition}
\newtheorem{corollary}[theorem]{Corollary}
\newtheorem{claim}[theorem]{Claim}

\theoremstyle{definition}
\newtheorem{definition}[theorem]{Definition}
\newtheorem{example}[theorem]{Example}
\newtheorem{exercise}[theorem]{Exercise}

\theoremstyle{remark}
\newtheorem*{remark}{Remark}
\newtheorem*{note}{Note}

%% ------------------------------------------------------------------
%%  Shorthand macros
%% ------------------------------------------------------------------
\newcommand{\R}{\mathbb{R}}
\newcommand{\N}{\mathbb{N}}
\newcommand{\Z}{\mathbb{Z}}
\newcommand{\Q}{\mathbb{Q}}
\newcommand{\eps}{\varepsilon}
\newcommand{\abs}[1]{\left\lvert #1 \right\rvert}
\newcommand{\norm}[1]{\left\lVert #1 \right\rVert}
\newcommand{\set}[1]{\left\{ #1 \right\}}
\newcommand{\ip}[1]{\left\langle #1 \right\rangle}
\DeclareMathOperator*{\argmax}{arg\,max}
\DeclareMathOperator*{\argmin}{arg\,min}

\title{Proofs \\[0.35em] \large Formal Logic}
\author{Ryan Gibbons \& Luke Cain \\ UChicago Political Science Math Camp}
\date{Math Camp 2026}

\begin{document}
\maketitle

\subsection*{Direct Democracy}

Contrapositive: every fair democracy has a tiebreaking rule. \\

Suppose that a country does not have a tiebreaking rule. Then, they either have elections without tiebreaks, or no elections at all. \\

If no elections, then clearly not a democracy.

If no tiebreaks, then someone must have outsized influence over the outcome (be decisive) -- but then, they cannot be a democracy. 


\subsection*{Median Voter Theorem}

Proof by contradiction: suppose not. Then, there exists some equilibrium whereby either $A$ or $B$ is playing a policy separate from the median voter's preferences. \\

Then, either $A$ is closer to $M$ than $B$, or $B$ is closer to $M$ than $A$, or $A$ and $B$ are some equal (non-zero) distance away from $M$. \\

If one is further than the other, the closer one will win -- so, the further one cannot be playing an optimal point. \\

If they are the same (non-zero) distance away, either one could win by moving just a little bit closer to $M$ -- so, neither of these players can be playing an optimal point. \\

Therefore, it cannot be the case that $A \neq M$ or $B \neq M$ in equilibrium. Therefore, $A = B = M$ in equilibrium.

\end{document}
